% Generated draft: requires mathematical and source review.
\documentclass[11pt]{article}
\usepackage[margin=25mm]{geometry}
\usepackage{fontspec}
\setmainfont{Latin Modern Roman}
\usepackage{amsmath,amssymb,mathtools,unicode-math}
\setmathfont{Latin Modern Math}
\usepackage{enumitem,needspace,xurl,hyperref}
\setlength{\parindent}{0pt}
\setlength{\parskip}{5pt}
\setlength{\emergencystretch}{3em}
\newcommand{\N}{\mathbb{N}}
\newcommand{\Z}{\mathbb{Z}}
\newcommand{\Q}{\mathbb{Q}}
\newcommand{\R}{\mathbb{R}}
\newcommand{\C}{\mathbb{C}}
\newcommand{\F}{\mathbb{F}}
\newcommand{\A}{\mathbb{A}}
\newcommand{\PP}{\mathbb{P}}
\newcommand{\E}{\mathbb{E}}
\newcommand{\eps}{\varepsilon}
\newcommand{\dd}{\,\mathrm{d}}
\begin{document}
\section*{UnsolvedMath ID: 30001157}
\textbf{Source problem code:} catalogue entry
\subsection{Definitions and mathematical statement}
Let \(\mathcal C\) denote the class of nonempty bounded open convex sets \(\Omega\subset\R^2\). Boundedness and nonempty interior are the conventional hypotheses used here to make the spectral optimization precise; they are implicit rather than explicit in the supplied statement. Write \(L(\Omega)=\mathcal H^1(\partial\Omega)\) for the perimeter. Define the first positive Neumann eigenvalue of the Laplacian by
\[
\mu_1(\Omega)=\inf_{\substack{u\in H^1(\Omega)\setminus\{0\}\\ \int_\Omega u\,\dd x=0}}
\frac{\int_\Omega |\nabla u|^2\,\dd x}{\int_\Omega |u|^2\,\dd x}.
\]
Thus the Neumann boundary condition is understood in the weak sense. Set \(F(\Omega)=\mu_1(\Omega)L(\Omega)^2\). Let \(Q=(0,1)^2\), and let \(T\) be the interior of the equilateral triangle with vertices \((0,0)\), \((1,0)\), and \((1/2,\sqrt3/2)\).

Determine
\[
\sup_{\Omega\in\mathcal C}F(\Omega)
\]
and all domains attaining it. In particular, is it true that
\[
F(\Omega)\leq F(Q)=F(T)\qquad\text{for every }\Omega\in\mathcal C,
\]
with equality if and only if \(\Omega\) is similar to \(Q\) or \(T\)? Here similarity allows translation, rotation, reflection, and positive uniform scaling.

The source additionally asks whether an underlying principle explains the equal values for the square and equilateral triangle. It does not specify a mathematical criterion for such an explanation, so this remains a qualitative part of the question.
\subsection{Short English statement}
Among bounded convex planar domains, determine the largest first positive Neumann eigenvalue multiplied by the square of the perimeter, and characterize every maximizing shape. The proposed maximizers are exactly squares and equilateral triangles, up to similarity. The source also asks for a conceptual explanation of their equal values; that request has no specified formal criterion.
\subsection{Sources}
\begin{itemize}
\item Supplied local UnsolvedMath registry record, permanent ID \(30001157\), title ``Perimeter-Normalized Neumann Eigenvalue Maximizers,'' source problem code \texttt{catalogue entry}, in the collection ``Oberwolfach Reports Open Problems.'' The record gives the proposed year as 2009 and names no proposer.
\item \emph{Low Eigenvalues of Laplace and Schrödinger Operators} (2009), Oberwolfach Reports, pp.\ 355--428, DOI \texttt{10.4171/owr/2009/06}: \url{https://doi.org/10.4171/owr/2009/06}. The supplied provenance describes Oberwolfach Reports as documenting mathematical workshops and their open-problem discussions.
\item Supplied literature reference: Henrot--Lemenant--Lucardesi, \url{https://doi.org/10.1090/tran/9074}.
\item Supplied literature reference: Laugesen--Siudeja, triangles, \url{https://doi.org/10.1063/1.3246834}.
\item Literature assessment dated exactly \texttt{2026-08-21}: the supplied curation labels the problem ``partially solved'' and reports that the equilateral triangle maximizes the functional within triangles; that the square and equilateral triangle are the only maximizers among convex domains with two axes of symmetry; and that the unrestricted convex-domain problem remains open. These are supplied dated curation claims, not independently verified mathematical or literature-status conclusions. The background explicitly requests independent verification of time-sensitive status and cited claims.
\item The registry labels the statement \texttt{corrected\_verified} and reports that its mathematical content was checked against the official Oberwolfach report, with notation, context, and question form normalized without changing the source-level claim. This verification description is supplied metadata; no independent source verification is asserted here.
\end{itemize}
\end{document}
